% ЗГЕНЕРОВАНО scripts/build_format_v2_chapter.py — не редагувати вручну.
% Джерело: books/ashtavakra-gita/poetychnyi/rozdil-10.md

\poemverse{10.1}{Покинув вороже бажання, \\ покинув користь, сповнену лиха, — \\ навіть дхарму, що є причиною обох. \\ Усюди зневажай.}
\poemverse{10.2}{Дивись на друзів, поле, гроші, \\ дім, дружину, спадок \\ і решту багатств — \\ як на сон і чари, \\ на три дні або пʼять.}
\poemverse{10.3}{Де б не виникла спрага — \\ знай: саме там і сансара. \\ Опираючись на зрілу безсторонність, \\ стань тим, у кого спрага минула, — \\ і будь щасливий.}
\poemverse{10.4}{Полон має природу самої лише спраги; \\ її знищення називають визволенням. \\ Від самої лише непривʼязаності до буття \\ приходить вдоволеність досягнутим — \\ знову й знову.}
\poemverse{10.5}{Ти один — свідомий, чистий; \\ світ же — неживий і несправжній. \\ Навіть незнання — аніщо. \\ То яке в тебе, при всьому тому, \\ бажання пізнавати?}
\poemverse{10.6}{Царство, сини, дружини, тіла й утіхи — \\ усе це вже гинуло для тебе, \\ привʼязаного, — \\ з народження в народження.}
\poemverse{10.7}{Досить уже й багатства, і бажання, \\ і навіть добрих учинків: \\ у цьому лісі сансари \\ твій розум так і не знав відпочинку.}
\poemverse{10.8}{Скільки ж за ці народження \\ зроблено тілом, розумом і мовою \\ діл, що дають біль і втому! \\ Тож нехай це припиниться хоч тепер.}
